\documentclass{article}
% ============================================================
% Bibliography and citations
% ============================================================

\usepackage{natbib}
\usepackage{csquotes}


% ============================================================
% Mathematics
% ============================================================

\usepackage{
    amsmath,
    amssymb,
    amsthm,
    amsfonts,
    mathrsfs,
    dsfont,
    bm,
    textgreek,
    bbm
}


% ============================================================
% Fonts
% ============================================================

\usepackage{newpxtext}
\usepackage{eulervm}


% ============================================================
% Page layout and document structure
% ============================================================

\usepackage[margin=1in, footskip=0.3in]{geometry}
\usepackage{setspace}
\usepackage[bf,small]{titlesec}
\usepackage{etoc}
\usepackage{chngcntr}
\usepackage{changepage}
\usepackage{afterpage}


% ============================================================
% Figures, captions, and graphics
% ============================================================

\usepackage{graphicx}
\usepackage[bf,skip=8pt]{caption}
\usepackage{subcaption}
\usepackage{tikz}

\graphicspath{{Figures/}}


% ============================================================
% Tables
% ============================================================

\usepackage{booktabs}
\usepackage{tabularx}
\usepackage{dcolumn}
\usepackage{makecell}
\usepackage{multirow}
\usepackage{longtable}
\usepackage[para]{threeparttable}
\usepackage{colortbl}
\usepackage{adjustbox}


% ============================================================
% Landscape and rotated content
% ============================================================

\usepackage{lscape}
\usepackage{pdflscape}
\usepackage{rotating}


% ============================================================
% Lists, comments, notes, and text utilities
% ============================================================

\usepackage{enumerate}
\usepackage{comment}
\usepackage{verbatim}
\usepackage{soul}
\usepackage[multiple,stable]{footmisc}
\usepackage[colorinlistoftodos]{todonotes}


% ============================================================
% Code listings
% ============================================================

\usepackage{listings}


% ============================================================
% Colors
% ============================================================

\usepackage{xcolor}

% General-purpose colors
\definecolor{blue}{rgb}{0,0.08,0.5}
\definecolor{red}{rgb}{0.6,0,0}
\definecolor{green}{rgb}{0,0.5,0.05}

% Code-listing colors
\definecolor{fgreen}{HTML}{698B27}
\definecolor{lavender}{HTML}{E6E6FA}
\definecolor{darkpink}{HTML}{C71585}
\definecolor{darkyellow}{HTML}{B8860B}


% ============================================================
% Hyperlinks and PDF metadata
% ============================================================

\usepackage[
    bookmarks,
    pdftitle={},
    pdfauthor={Banks},
    colorlinks=true,
    linkcolor=black,
    citecolor=black,
    urlcolor=black
]{hyperref}


% ============================================================
% Theorem environments
% ============================================================

\newtheorem{prediction}{Prediction}
\newtheorem{assumption}{Assumption}
\newtheorem{proposition}{Proposition}
\newtheorem{corollary}{Corollary}
\newtheorem{definition}{Definition}
\newtheorem{lemma}{Lemma}

% Proposition numbering resets by section
\newtheorem{prop}{Proposition}[section]


% ============================================================
% Mathematical commands
% ============================================================

% Formatting
\newcommand{\no}{\noindent}
\newcommand{\mb}{\mathbf}

% Expectation, probability, and covariance
\newcommand{\E}{\mathbb{E}}
\newcommand{\R}{\mathbb{R}}
\newcommand{\pr}{\mathrm{Pr}}
\newcommand{\cov}{\mathrm{Cov}}

% Support
\newcommand{\supp}[1]{\mathrm{Supp}\left(#1\right)}

% Partial derivatives
\newcommand{\pd}[2]{\frac{\partial #1}{\partial #2}}
\newcommand{\pdd}[2]{\frac{\partial^2 #1}{\partial #2^2}}

% Indicator function
\newcommand{\ind}[1]{\mathbbm{1}\{#1\}}

% Vertical spacing helper for large fractions
\newcommand{\vh}{
    \vphantom{
        \frac{
            \frac{\pd{a}{a}}{\pd{a}{a}}
        }{
            \frac{\pd{a}{a}}{\pd{a}{a}}
        }
    }
}


% ============================================================
% Custom environments
% ============================================================

% Wide equation extending 2 cm into each margin
\newenvironment{widerequation}
{
    \begin{adjustwidth}{-2cm}{-2cm}
    \begin{equation*}
}
{
    \end{equation*}
    \end{adjustwidth}
}

% Solution environment
\newenvironment{solution}
{
    \par\noindent\textbf{\textit{Solution.}}
}
{
    \hfill\rule{1.1ex}{1.1ex}
}


% ============================================================
% Code-listing style
% ============================================================

\lstset{
    basicstyle=\ttfamily\small\color{fgreen},
    keywordstyle=\color{darkpink},
    commentstyle=\color{gray},
    stringstyle=\color{darkyellow},
    showstringspaces=false,
    breaklines=true,
    frame=single,
    backgroundcolor=\color{lavender}
}


% ============================================================
% Spacing and section formatting
% ============================================================

% ECMA-style line spacing
\linespread{1.5}

\titlespacing*{\paragraph}
    {0pt}
    {.25\baselineskip}
    {1em}

\titlespacing*{\section}
    {0pt}
    {.6\baselineskip}
    {.3\baselineskip}

\titlespacing*{\subsection}
    {0pt}
    {.6\baselineskip}
    {.3\baselineskip}

\titlespacing*{\subsubsection}
    {0pt}
    {.6\baselineskip}
    {.3\baselineskip}

% \setlength{\parskip}{5pt}

\setlength{\footnotesep}{.2in}


% ============================================================
% Subsubsection formatting
% ============================================================

\titleformat{\subsubsection}[runin]
    {\normalfont\itshape\bfseries}
    {\thesubsubsection}
    {1em}
    {}
    [.]


% ============================================================
% Date and title-page utilities
% ============================================================

\newcommand{\monthyear}{%
    \ifcase\month\or
        January\or
        February\or
        March\or
        April\or
        May\or
        June\or
        July\or
        August\or
        September\or
        October\or
        November\or
        December
    \fi
    \ \number\year
}


\begin{document}
    \begin{center}
        \scalebox{1.3}{Economics 880 - PS4}\\
        Aaron Banks \& Arthur Rodrigues
    \end{center}

\section*{Question 1}

The moment vector $m_3(x_t)$ contains the quantities used to
construct the mean, variance, and first-order autocovariance of
the time series $x \equiv \{x_t\}$. We have
\[
    m_3(x_t)
    =
    \begin{bmatrix}
        x_t\\
        (x_t-\bar{x})^2\\
        (x_t-\bar{x})(x_{t-1}-\bar{x})
    \end{bmatrix}.
\]

We will consider the moments of the stationary distribution of the
series since we are asked to compute the asymptotic moments.
We have
\begin{align*}
    \mathbb{E}[x_t]
    &= \rho_0 \mathbb{E}[x_{t-1}]
    + \mathbb{E}[\varepsilon_t]\\
    &= \rho_0 \mathbb{E}[x_{t-1}].
\end{align*}
Under stationarity, $\mathbb{E}[x_t]=\mu$ for all $t$, hence
\[
    \mu = \rho_0\mu
    \qquad \Rightarrow \qquad
    \mu(1-\rho_0)=0
    \qquad \Rightarrow \qquad
    \mu=0,
\]
since $\rho_0=0.5$.

For the asymptotic variance, we have
\begin{align*}
    \mathbb{E}[x_t^2]
    &= \rho_0^2\mathbb{E}[x_{t-1}^2]
    + 2\rho_0\mathbb{E}[x_{t-1}\varepsilon_t]
    + \mathbb{E}[\varepsilon_t^2]\\
    &= \rho_0^2\mathbb{E}[x_{t-1}^2]+\sigma_0^2,
\end{align*}
since $\mathbb{E}[x_{t-1}\varepsilon_t]=0$.
Under stationarity,
$\mathbb{E}[x_{t-1}^2]=\mathbb{E}[x_t^2]=v$, so
\[
    v=\rho_0^2v+\sigma_0^2
    \qquad \Rightarrow \qquad
    v=\frac{\sigma_0^2}{1-\rho_0^2}.
\]

For the first-order autocovariance, we have
\begin{align*}
    \mathbb{E}[x_tx_{t-1}]
    &= \mathbb{E}[
        (\rho_0x_{t-1}+\varepsilon_t)x_{t-1}
    ]\\
    &= \rho_0\mathbb{E}[x_{t-1}^2]
    + \mathbb{E}[\varepsilon_tx_{t-1}]\\
    &= \rho_0\mathbb{E}[x_{t-1}^2]\\
    &= \frac{\rho_0\sigma_0^2}{1-\rho_0^2},
\end{align*}
where the third equality follows from the independence of
$\varepsilon_t$ and $x_{t-1}$, and the fourth from stationarity.

Thus, for the parameter vector
$b=(\rho,\sigma^2)^{\mathsf T}$, the population moment vector is
\[
    \mu_3(b)
    =
    \begin{bmatrix}
        0\\[6pt]
        \displaystyle\frac{\sigma^2}{1-\rho^2}\\[10pt]
        \displaystyle\frac{\rho\sigma^2}{1-\rho^2}
    \end{bmatrix}.
\]

Evaluating at the true parameters,
$b_0=(\rho_0,\sigma_0^2)^{\mathsf T}$, gives
\[
    \mu_3(b_0)
    =
    \begin{bmatrix}
        0\\[6pt]
        \displaystyle\frac{\sigma_0^2}{1-\rho_0^2}\\[10pt]
        \displaystyle\frac{\rho_0\sigma_0^2}{1-\rho_0^2}
    \end{bmatrix}
    =
    \begin{bmatrix}
        0\\
        4/3\\
        2/3
    \end{bmatrix}.
\]

We next compute the Jacobian of the population moment conditions.
Recall that
\[
    g(b)=\mu_3(b_0)-\mu_3(b),
\]
so
\[
    \nabla_b g(b)
    =
    -\frac{\partial\mu_3(b)}{\partial b^{\mathsf T}}.
\]

Differentiating with respect to $\rho$ and $\sigma^2$, we obtain
\[
    \nabla_b g(b)
    =
    -
    \begin{bmatrix}
        0 & 0\\[6pt]
        \displaystyle\frac{2\rho\sigma^2}{(1-\rho^2)^2}
        &
        \displaystyle\frac{1}{1-\rho^2}\\[10pt]
        \displaystyle\frac{\sigma^2(1+\rho^2)}{(1-\rho^2)^2}
        &
        \displaystyle\frac{\rho}{1-\rho^2}
    \end{bmatrix}.
\]

Evaluating at $b_0$, we have
\[
    G \equiv \nabla_b g(b_0)
    =
    -
    \begin{bmatrix}
        0 & 0\\
        16/9 & 4/3\\
        20/9 & 2/3
    \end{bmatrix}.
\]

The Jacobian has full column rank since
\[
    \det
    \begin{bmatrix}
        16/9 & 4/3\\
        20/9 & 2/3
    \end{bmatrix}
    =
    -\frac{16}{9}\neq 0.
\]
Thus,
\[
    \operatorname{rank}(G)=2,
\]
and the two structural parameters are locally identified.

We next compute the AGS sensitivity matrix, defined as
\[
    \Lambda
    =
    -(G^{\mathsf T}WG)^{-1}G^{\mathsf T}W.
\]

Using $W=I_3$, we have
\[
    \Lambda
    =
    -(G^{\mathsf T}G)^{-1}G^{\mathsf T}.
\]

Substituting the Jacobian and simplifying gives
\[
    \Lambda
    =
    \begin{bmatrix}
        0 & -3/8 & 3/4\\
        0 & 5/4 & -1
    \end{bmatrix}.
\]

The first column of $\Lambda$ is zero, indicating that the
sample mean does not affect either parameter estimate when
$W=I_3$. This follows from the fact that the population mean
is zero regardless of the structural parameters.

The variance and first-order autocovariance are informative
for both parameters. In particular, the first row of $\Lambda$
indicates that an increase in the empirical variance decreases
the estimate of $\rho$, holding the autocovariance fixed,
while an increase in the empirical autocovariance increases
the estimate of $\rho$, holding the variance fixed.

Similarly, the second row indicates that an increase in
the empirical variance increases the estimate of $\sigma^2$,
holding the autocovariance fixed, while an increase in the
empirical autocovariance decreases the estimate of $\sigma^2$,
holding the variance fixed.

Thus, the mean provides no identifying information, while
the variance and autocovariance jointly identify the two
structural parameters.


\end{document}